\documentclass[11pt,a4paper]{article}
\usepackage[utf8]{inputenc}
\usepackage[T1]{fontenc}
\usepackage{amsmath,amssymb,amsthm,mathtools}
\usepackage{geometry}
\usepackage{enumitem}
\usepackage{fancyhdr}
\usepackage{titlesec}
\usepackage{microtype}

\geometry{margin=1in}

\pagestyle{fancy}
\fancyhf{}
\fancyhead[C]{\small Cayley, Collected Mathematical Papers, Vol.~IV}
\fancyfoot[C]{\thepage}
\renewcommand{\headrulewidth}{0.4pt}

\titleformat{\section}{\large\bfseries}{\thesection.}{0.5em}{}
\titleformat{\subsection}{\normalsize\bfseries}{\thesubsection.}{0.5em}{}

\newtheorem*{theorem}{Theorem}
\newcommand{\paper}[2]{\section[#1]{#1\hfill\textnormal{[#2]}}}

\begin{document}

\begin{center}
{\Large\bfseries The Collected Mathematical Papers of Arthur Cayley}\\[6pt]
{\large Vol.~IV, Pages 501--550}\\[12pt]
{\small Transcribed from the original publication}
\end{center}

\vspace{1em}
\hrule
\vspace{1.5em}

\paper{On the Construction of the Ninth Point of Intersection of the Cubics which pass through Eight Given Points}{295}

{\small
[\textit{From the Quarterly Journal of Pure and Applied Mathematics, vol.~v. (1862), pp. 222--233.}]
\par}
\vspace{0.5em}

\setcounter{page}{501}

anharmonic ratio $9(4, 6, 7, 8)$ is equal to that of the points $(4, 6, Z, L)$ on the line $46Z_1$, that is
\[
9 (4, 6, 7, 8) = (4, 6, Z, L),
\]
which is
\[
= (4, M, 7, Q),
\]
since the lines $6M$, $Zl$, $LQ$ meet in the point $K$; but by the construction
\[
(4, M, 7, Q) = 1235\,(4, 6, 7, 8),
\]
that is, we have
\[
9 (4, 6, 7, 8) = 1235\,(4, 6, 7, 8);
\]
and, secondly, the anharmonic ratio $9(5, 6, 7, 8)$ is equal to that of the points $(5, 6, K, Y)$ on the line $56KY$, that is
\[
9(5, 6, 7, 8) = (5, 6, K, Y),
\]
which is
\[
= (5, N, R, 8),
\]
since the lines $6N$, $KR$, $Y8$ meet in the same point $L$; but by the construction
\[
(5, N, R, 8) = 1234\,(5, 6, 7, 8),
\]
that is, we have
\[
9(5, 6, 7, 8) = 1234\,(5, 6, 7, 8),
\]
so that the point $9$ satisfies the required conditions.

It has been already remarked that the point $x$ in M.~Chasles' theorem for the construction of the cubic through nine points is determined by precisely similar conditions to those which determine the ninth intersection of the two cubics; that is, the foregoing construction by the ruler alone is applicable to the determination of the point $x$; and when this is once obtained, the remainder of the construction, giving the points of the cubic through the nine given points, can obviously be performed by the ruler alone. The construction for the cubic through nine points gives implicitly the relation between ten points of the cubic and such relation is accordingly expressed by the equation
\[
x(1, 2, 3, 4, 5, 10) = 6789\,(1, 2, 3, 4, 5, 10),
\]
which is one out of 210 similar forms. But it is possible that some more convenient form of the relation between the ten points may yet be found.

I proceed to further develope the analytical theory. Writing for convenience $\omega$ in the place of $10$, we have
\begin{align*}
x(1, 2, 3, 4, 5, 6) &= 789\omega\,(1, 2, 3, 4, 5, 6),\\
\shortintertext{or, what is the same thing,}
x(1, 2, 3, 4) &= 789\omega\,(1, 2, 3, 4),\\
x(1, 2, 3, 5) &= 789\omega\,(1, 2, 3, 5),\\
x(1, 2, 3, 6) &= 789\omega\,(1, 2, 3, 6),
\end{align*}

\newpage
\setcounter{page}{502}

which belong to three conics each of them passing through $1$, $2$, $3$, and which must have a remaining fourth point of intersection.

The equation of the first conic is
\[
012 \cdot 034 = -\lambda \cdot 014 \cdot 023,
\]
where
\begin{align*}
\lambda &= 789\omega 12 \cdot 789\omega 34,\\
\mu &= 789\omega 14 \cdot 789\omega 23,
\end{align*}
and thence, in virtue of an identical equation already referred to,
\[
-\lambda - \mu = 789\omega 13 \cdot 789\omega 24.
\]
But we have identically
\[
123 \cdot 0pq = 1pq \cdot 023 + 2pq \cdot 031 + 3pq \cdot 012,
\]
and thence in particular
\begin{align*}
123 \cdot 034 &= 134 \cdot 023 + 234 \cdot 031,\\
123 \cdot 014 &= 214 \cdot 031 + 314 \cdot 012,
\end{align*}
and the equation of the conic may therefore be written
\[
012 (134 \cdot 023 + 234 \cdot 031)\mu - 023 (214 \cdot 031 + 314 \cdot 012)\lambda = 0,
\]
that is
\[
031 \cdot 012 \cdot 234 \cdot \mu + 012 \cdot 023 \cdot 314(-\lambda-\mu) + 023 \cdot 031 \cdot 124 \cdot \lambda = 0,
\]
or, substituting for $\mu$, $-\lambda-\mu$, $\lambda$, their values, this is
\begin{multline*}
031 \cdot 012 \cdot 234 \cdot 789\omega 14 \cdot 789\omega 23\\
+ 012 \cdot 023 \cdot 314 \cdot 789\omega 13 \cdot 789\omega 42\\
+ 023 \cdot 031 \cdot 124 \cdot 789\omega 12 \cdot 789\omega 34 = 0,
\end{multline*}
or, what is the same thing,
\[
\frac{234 \cdot 789\omega 14 \cdot 789\omega 23}{023}
+ \frac{314 \cdot 789\omega 13 \cdot 789\omega 24}{031}
+ \frac{124 \cdot 789\omega 12 \cdot 789\omega 34}{012} = 0,
\]
or, making a slight change of form, the equation of the conic is
\[
\frac{423 \cdot 789\omega 41 \cdot 789\omega 23}{023}
+ \frac{431 \cdot 789\omega 42 \cdot 789\omega 31}{031}
+ \frac{412 \cdot 789\omega 43 \cdot 789\omega 12}{012} = 0.
\]
The equations of the other two conics are deduced by writing successively $5$ and $6$ in the place of $4$; and the condition in order that the conics may have a remaining fourth point of intersection is
\[
\begin{vmatrix}
423 \cdot 789\omega 41, & 431 \cdot 789\omega 42, & 412 \cdot 789\omega 43\\
523 \cdot 789\omega 51, & 531 \cdot 789\omega 52, & 512 \cdot 789\omega 53\\
623 \cdot 789\omega 61, & 631 \cdot 789\omega 62, & 612 \cdot 789\omega 63
\end{vmatrix}
= 0.
\]

\newpage
\setcounter{page}{503}

This equation, say $\Omega = 0$, expresses the relation between the coordinates of the ten points $1$, $2$, $3$, $4$, $5$, $6$, $7$, $8$, $9$, $\omega$, of the cubic. Hence if $123456789\omega$ denote the determinant formed with the ten quantities
\[
(x^3, y^2z, z^2x, x^2y, y^3, z^3, x^2z, y^2x, z^2y, xyz)
\]
for the ten points respectively, $123456789\omega$ must be a factor of $\Omega$, and a little consideration shows that the other factor which is of the order one as regards the coordinates of each of the points $1$, $2$, $3$, and of the order three as regards the coordinates of each of the points $7$, $8$, $9$, $\omega$, must be of the form $123 \cdot 789 \cdot 78\omega \cdot 79\omega \cdot 89\omega$. We must therefore have
\[
\Omega = \epsilon \cdot 123 \cdot 789 \cdot 78\omega \cdot 79\omega \cdot 89\omega \cdot 123456789\omega,
\]
where the merely numerical factor $\epsilon$ is, I believe, equal to $+1$ or else to $-1$.

In order to verify the factor $123 \cdot 789 \cdot 78\omega \cdot 79\omega \cdot 89\omega$, observing that the points $7$, $8$, $9$, $\omega$ enter symmetrically, it will be sufficient to show that $123$, $789$ are each of them factors of $\Omega$, or, what is the same thing, that if $123 = 0$, or if $789 = 0$, then in either case $\Omega = 0$.

First, if $123 = 0$, we may write
\begin{align*}
x_3 &= \lambda x_1 + \mu x_2,\\
y_3 &= \lambda y_1 + \mu y_2,\\
z_3 &= \lambda z_1 + \mu z_2,
\end{align*}
equations which give
\[
423 = \lambda \cdot 421,\quad 431 = \mu \cdot 421,\quad \&\mathrm{c.},
\]
and the equation $\Omega = 0$, thus becomes
\[
\begin{vmatrix}
789\omega 41, & 789\omega 42, & 789\omega 43\\
789\omega 51, & 789\omega 52, & 789\omega 53\\
789\omega 61, & 789\omega 62, & 789\omega 63
\end{vmatrix}
= 0,
\]
or, what is the same thing,
\[
\begin{vmatrix}
789\omega 14, & 789\omega 24, & 789\omega 34\\
789\omega 15, & 789\omega 25, & 789\omega 35\\
789\omega 16, & 789\omega 26, & 789\omega 36
\end{vmatrix}
= 0.
\]
Now if the terms in the same column are multiplied by $789\omega 56$, $789\omega 64$, $789\omega 45$ respectively and added, then for the first column the sum is
\[
789\omega 14 \cdot 789\omega 56 + 789\omega 15 \cdot 789\omega 64 + 789\omega 16 \cdot 789\omega 45,
\]
which is $= 0$, and the sums for the second and third columns are each $= 0$ in the same manner: wherefore the determinant vanishes as it should do.

\newpage
\setcounter{page}{504}

Next, if $789 = 0$, we have identically
\[
789\omega 41 = 789 \cdot 741 \cdot \omega 81 \cdot \omega 94 + 781 \cdot 794 \cdot \omega 89 \cdot \omega 41,
\]
which when $789 = 0$ gives
\[
789\omega 41 = 781 \cdot 794 \cdot \omega 89 \cdot \omega 41,
\]
and in like manner,
\begin{align*}
789\omega 42 &= 782 \cdot 794 \cdot \omega 89 \cdot \omega 42,\\
789\omega 43 &= 783 \cdot 794 \cdot \omega 89 \cdot \omega 43,
\end{align*}
$= 0$, in which three equations $4$ may be changed into $5$ and $6$ successively. The equation $\Omega = 0$ thus becomes
\[
\begin{vmatrix}
423 \cdot \omega 41, & 431 \cdot \omega 42, & 412 \cdot \omega 43\\
523 \cdot \omega 51, & 531 \cdot \omega 52, & 512 \cdot \omega 53\\
623 \cdot \omega 61, & 631 \cdot \omega 62, & 612 \cdot \omega 63
\end{vmatrix}
\]
or, what is the same thing,
\[
\begin{vmatrix}
423 \cdot 4\omega 1, & 421 \cdot 4\omega 3, & 42\omega \cdot 431\\
523 \cdot 5\omega 1, & 521 \cdot 5\omega 3, & 52\omega \cdot 531\\
623 \cdot 6\omega 1, & 621 \cdot 6\omega 3, & 62\omega \cdot 631
\end{vmatrix}
= 0,
\]
and since the sum of the terms in each line of the determinant is $=0$, the determinant is as it should be $= 0$.

The foregoing equation
\[
\begin{vmatrix}
412 \cdot 789\omega 43, & 423 \cdot 789\omega 41, & 431 \cdot 789\omega 42\\
512 \cdot 789\omega 53, & 523 \cdot 789\omega 51, & 531 \cdot 789\omega 52\\
612 \cdot 789\omega 63, & 623 \cdot 789\omega 61, & 631 \cdot 789\omega 62
\end{vmatrix}
= \epsilon \cdot 123 \cdot 789 \cdot 78\omega \cdot 79\omega \cdot 89\omega \cdot 123456789\omega,
\]
(since $412$, $789\omega 43$, \&c.\ are interpretable functions of the coordinates) affords a geometrical interpretation of the equation
\[
123456789\omega = 0
\]
between the coordinates of the ten points of the cubic; but it would be more satisfactory if a similar identical equation could be found, having on the right-hand side the function $123456789\omega$ without the irrelevant factor $123 \cdot 789 \cdot 78\omega \cdot 79\omega \cdot 89\omega$.

\begin{flushright}
2, Stone Buildings, W.C., 6th March, 1862.
\end{flushright}

\vspace{1.5em}
\hrule
\vspace{1.5em}

\paper{On the Conics which pass through the Four Foci of a Given Conic}{296}

{\small
[\textit{From the Quarterly Journal of Pure and Applied Mathematics, vol.~v. (1862), pp. 275--280.}]
\par}
\vspace{0.5em}

\setcounter{page}{505}

The foci of a conic are the points of intersection of the tangents through the circular points at infinity; the pair of tangents through each of the circular points at infinity is a conic through the four foci; and we have thus two conics $P = 0$, $Q = 0$ passing through the four foci; the equation of any other conic through the four foci is of course $P + \lambda Q = 0$; and in particular if $\lambda$ be suitably determined this equation gives the axes of the conic.

I was led to develope the solution, in seeking to obtain the elegant formulae given in Mr P.~J.~Hensley's paper ``Determination of the foci of the conic section expressed by trilinear coordinates,'' \textit{Journal}, t.~v., pp.~177---183, (March, 1862).

I take the coordinates to be proportionate to the perpendicular distances of the point from the sides of the fundamental triangle, each distance divided by the perpendicular distance of the side from the opposite angle. This being so, the equation of the line infinity is
\[
x + y + z = 0,
\]
and, $a$, $b$, $c$ denoting the sides of the fundamental triangle, the equation of the circle circumscribed about the triangle is
\[
\frac{a^2}{x} + \frac{b^2}{y} + \frac{c^2}{z} = 0.
\]
The foregoing two equations determine the circular points at infinity; and if $(x_1, y_1, z_1)$ are the coordinates, there is no difficulty in obtaining the system of values
\begin{align*}
x_1 &: y_1 : z_1\\
&= a(\cos C - i \sin C) : b(\cos B + i \sin B) : -c\\
&= b(\cos C + i \sin C) : -a : c(\cos A - i \sin A)\\
&= -b : a(\cos B + i \sin B) : c(\cos A + i \sin A),
\end{align*}

\newpage
\setcounter{page}{506}

where as usual $i = \sqrt{-1}$, and where $A$, $B$, $C$ denote the angles of the triangle, so that the cosines and sines of these angles denote given functions of $a$, $b$, $c$. The coordinates $(x_2, y_2, z_2)$ of the other circular point at infinity are of course obtained by merely writing $-i$ for $i$. We find also
\begin{multline*}
x_1x_2 : y_1y_2 : z_1z_2 : y_1z_2 + y_2z_1 : z_1x_2 + z_2x_1 : x_1y_2 + x_2y_1\\
= -a^2 : -b^2 : -c^2 : b^2 + c^2 - a^2 : c^2 + a^2 - b^2 : a^2 + b^2 - c^2,
\end{multline*}
which are the formula chiefly made use of in the sequel.

Suppose now that the equation of the conic is
\[
U = (a, b, c, f, g, h \between x, y, z)^2 = 0;
\]
then putting for a moment
\begin{align*}
U_1 &= (a, b, c, f, g, h \between x_1, y_1, z_1)^2,\\
V_1 &= (a, b, c, f, g, h \between x, y, z \between x_1, y_1, z_1),
\end{align*}
and the like as regards $U_2$ and $V_2$; the tangents from the points $(x_1, y_1, z_1)$ and $(x_2, y_2, z_2)$ respectively are
\begin{align*}
UU_1 - V_1^2 &= 0,\\
UU_2 - V_2^2 &= 0,
\end{align*}

\newpage
\setcounter{page}{507}

and this being so, I combine the two equations as follows:
\begin{align*}
x_1{}^2 (a, \ldots)(x_1, y_1, z_1)^2 + x_2{}^2 (a, \ldots)(x_2, y_2, z_2)^2 &= 0,\\
y_1{}^2 \,\,\,\, ,, \quad\quad + y_2{}^2 \,\,\,\, ,, \quad\quad &= 0,\\
z_1{}^2 (a, \ldots)(x_1, y_1, z_1)^2 + z_2{}^2 (a, \ldots)(x_2, y_2, z_2)^2 &= 0,\\
y_1z_1 \,\,\,\, ,, \quad\quad + y_2z_2 \,\,\,\, ,, \quad\quad &= 0,\\
z_1x_1 \,\,\,\, ,, \quad\quad + z_2x_2 \,\,\,\, ,, \quad\quad &= 0,\\
x_1y_1 \,\,\,\, ,, \quad\quad + x_2y_2 \,\,\,\, ,, \quad\quad &= 0,
\end{align*}
any one of which is the equation of a conic passing through the four foci; the current coordinates being always $(x, y, z)$.

The first of these equations is
\begin{multline*}
a(-2x_1{}^2x_2{}^2) + b(x_2{}^2y_1{}^2 + x_1{}^2y_2{}^2) + c(x_2{}^2z_1{}^2 + x_1{}^2z_2{}^2)\\
+ 2f(x_2{}^2y_1z_1 + x_1{}^2y_2z_2) + 2g(x_1x_2{}^2z_1 + x_1{}^2x_2z_2) + 2h(x_1x_2{}^2y_1 + x_1{}^2x_2y_2) = 0,
\end{multline*}
where the quantities multiplied by $a$, $b$, \&c.\ are all of them easily expressible in terms of $x_1x_2$, $y_1y_2$, $z_1z_2$, $y_1z_2 + y_2z_1$, $z_1x_2 + z_2x_1$, $x_1y_2 + x_2y_1$, which are respectively proportional to given functions of $(a, b, c)$; and replacing for $a$, $b$, \&c.\ their values, the equation is
\begin{multline*}
(-2c^2x^2 + 2b^2y^2 + 2a^2z^2 + 2(-c^2-a^2+b^2)yz + 2(-b^2+a^2+c^2)zx + 2(a^2+b^2-c^2)xy) \cdot 2a^2\\
- a^2(b^2+c^2-a^2)\cdot(\ldots) + \ldots
\end{multline*}

\newpage
\setcounter{page}{508}

or, what is the same thing, the equation is
\[
\Omega\{(2b + 2c + 2f + 2g + 2h)x^2 - 2(2b + 2f + 2h)x(x+y+z) + b(x+y+z)^2\} + 2a^2\Theta = 0,
\]
or putting for shortness
\begin{align*}
2b + 2c + 2f + 2g + 2h &= 2l,\\
2b + 2f + 2h &= 2m,\\
2c + 2f + 2g &= 2n,\\
2g + 2h + 2f &= 2p,
\end{align*}
so that in fact
\begin{align*}
l &= b+c+f+g+h,\\
m &= b+f+h,\\
n &= c+f+g,\\
p &= g+h+f,
\end{align*}
the equation is
\[
D\{lx^2 - 2mx(x+y+z) + b(x+y+z)^2\} + 2a^2\Theta = 0.
\]

The equation with $y_1y_2$, $y_1z_2$ is in a similar manner found to be
\[
D\{(2b+2c+2f+2g+2h)y^2 - \{(2c+2f+2g)y + (2b+2f+2h)z\}(x+y+z) + c(x+y+z)^2\} - a^2\Theta = 0,
\]
or, what is the same thing,
\[
D\{ly^2 - 2ny(x+y+z) + c(x+y+z)^2\} - a^2\Theta = 0,
\]
or, finally,
\[
D\{lz^2 - 2pz(x+y+z) + a(x+y+z)^2\} + 2c^2\Theta = 0.
\]

Hence the entire system of equations is
\begin{align*}
D\{lx^2 - 2mx(x+y+z) + b(x+y+z)^2\} + 2a^2\Theta &= 0,\\
D\{ly^2 - 2ny(x+y+z) + c(x+y+z)^2\} + 2b^2\Theta &= 0,\\
D\{lz^2 - 2pz(x+y+z) + a(x+y+z)^2\} + 2c^2\Theta &= 0,\\
D\{2fyz - (ny + mz)(x+y+z) + 2l(x+y+z)^2\} - (b^2+c^2-a^2)\Theta &= 0,\\
D\{2gzx - (pz + mx)(x+y+z) + 2m(x+y+z)^2\} - (c^2+a^2-b^2)\Theta &= 0,\\
D\{2hxy - (lx + py)(x+y+z) + 2n(x+y+z)^2\} - (a^2+b^2-c^2)\Theta &= 0,
\end{align*}
which are the equations of six conics, each of them passing through the four foci.

From the first three of these, we have
\begin{multline*}
\{lx^2 - 2mx(x+y+z) + b(x+y+z)^2\}\\
= \{ly^2 - 2ny(x+y+z) + c(x+y+z)^2\}\\
= \{lz^2 - 2pz(x+y+z) + a(x+y+z)^2\},
\end{multline*}

\newpage
\setcounter{page}{509}

which, allowing for the difference of notation, are Mr Hensley's equations: it appears by his investigation that their geometrical signification is as follows; viz. if for shortness we denote the equations by
\[
A = 0,\quad B = 0,\quad C = 0,
\]
then if we consider the tangents parallel to the $x$-side of the fundamental triangle, and the tangents parallel to the $y$-side of the fundamental triangle, the equation $A = B$ is the locus of a point such that the feet of the perpendiculars let fall from it on the four tangents lie in a circle. And similarly for the equations $A = C$, $B = C$.

If we multiply the six equations by $1$, $1$, $1$, $2$, $2$, $2$ and add, we obtain the identical equation $0 = 0$; if we multiply them by $a$, $b$, $c$, $2f$, $2g$, $2h$ and add, then after some easy reductions, we obtain for the equation of a new conic passing through the four foci
\[
D\{KU + K(x+y+z)^2\} + 2S\Theta = 0,
\]
where
\[
U = (a, b, c, f, g, h \between x, y, z)^2,
\]
$K$ is the discriminant $abc - af^2 - bg^2 - ch^2 + 2fgh$, and
\[
S = (b^2+c^2-a^2)f + (c^2+a^2-b^2)g + (a^2+b^2-c^2)h,
\]
or, what is the same thing,
\[
S = (f+a-h-g)a^2 + (g-h+b-f)b^2 + (h-g-f+c)c^2.
\]
It would be interesting to ascertain the geometrical signification of the six conics and of the last-mentioned new conic.

\begin{flushright}
2, Stone Buildings, W.C., 13th March, 1862.
\end{flushright}

\vspace{1.5em}
\hrule
\vspace{1.5em}

\paper{On some Formulae relating to the Distances of a Point from the Vertices of a Triangle, and to the Problem of Tactions}{297}

{\small
[\textit{From the Quarterly Journal of Pure and Applied Mathematics, vol.~v. (1862), pp. 381--384.}]
\par}
\vspace{0.5em}

\setcounter{page}{511}

if $D$ denote the function in $\{\quad\}$ with the signs reversed. The function $D$ may be expressed in the form
\begin{multline*}
+ (a^2p + b^2c^2)(f^2 - g^2 - h^2)\\
+ (b^2c^2q + c^2a^2r)(g^2 - h^2 - f^2)\\
+ (c^2a^2r + a^2b^2p)(h^2 - f^2 - g^2),
\end{multline*}
and also in the form
\begin{multline*}
\square U^2 + (f+g+h)V,
\end{multline*}
if for shortness
\begin{align*}
\square &= U^2 + (f+g+h)V,\\
U &= a^2f + b^2g + c^2h + fgh,\\
V &= (a^2p + b^2c^2)(f-g-h)\\
&\quad + (b^2q + c^2a^2)(g-h-f)\\
&\quad + (c^2r + a^2b^2)(h-f-g);
\end{align*}
and it may be remarked that since $D$ is an even function of $f$, $g$, $h$, we may in this last formula change at pleasure the signs of these quantities; we thus obtain in all four similar forms of the function $D$.

It is clear that considering a triangle, and any point in the plane of the triangle, $f$, $g$, $h$ may be taken to denote the sides of the triangle, and $a$, $b$, $c$ the distances of the point from the vertices: and the equation $D = 0$ is the relation connecting the sides and distances.

The equation $f+g+h=0$ denotes that the vertices are in lined, and when this equation is satisfied we have
\[
U = a^2f + b^2g + c^2h + fgh = 0,
\]
which is in fact, as it is easy to see, the relation connecting the distances of a point from any three points in lined.

For $a$, $b$, $c$ write $a+x$, $b+x$, $c+x$; $x$ will be the radius of a circle touching the circles, radii $a$, $b$, $c$, described about the vertices as centres. The equation $D = 0$ becomes after all reductions
\begin{multline*}
U^2 - (f+g+h)V\\
+ x[4U(af+bg+ch) - 2(f+g+h)\{(a^2p+bc(b+c))(f-g-h)\\
+ (b^2q+ca(c+a))(g-h-f) + (c^2r+ab(a+b))(h-f-g)\}]\\
+ x^2[f^2\{-4a^2+6a(b+c)-6bc\} + g^2\{-4b^2+6b(c+a)-6ca\}\\
+ h^2\{-4c^2+6c(a+b)-6ab\}] = 0,
\end{multline*}
which is a quadratic equation only: the two circles thus obtained are those which touch the given circles all three externally or all three internally. But by changing in every possible manner the signs of $a$, $b$, $c$ we obtain in all four equations giving the eight tangent circles. It may be noticed that if as before $f+g+h=0$, $U=0$,

\newpage
\setcounter{page}{512}

then not only the constant term vanishes, but the coefficient of $x$ also vanishes or the equation becomes simply $x^2 = 0$.

In particular, suppose $f=b+c$, $g=c+a$, $h=a+b$; developing this \textit{de novo}, and putting for shortness
\begin{align*}
a+b+c &= p,\\
bc+ca+ab &= q,\\
abc &= r,
\end{align*}
we find
\begin{align*}
U &= 2\{pa^2 + 4qa + (2q^2 - 2r)\},\\
V &= 2\{px^2 + 4qax + (2pq + 12r)a^2 + 4q^2a + pq^2 - 4qr\},
\end{align*}
and then the equation $D = U^2 - 2pV = 0$ gives
\begin{align*}
D &= (pa^2 + 2qa + pq - 2r)^2 - p\{pa^2 + 4qa + (2pq+12r)a^2 + 4q^2a + pq^2 - 4qr\}\\
&= 4\{(q^2 - 4pr)a^2 - 2qra + r^2\},
\end{align*}
so that we have
\[
\tfrac{1}{4}D = (q^2 - 4pr)a^2 - 2qra + r^2 = (qa-r)^2 - 4pra^2 = 0,
\]
and thence
\[
qa - r = \pm a\sqrt{pr},\quad \text{or } x = \frac{r}{q \mp \sqrt{pr}},
\]
which gives the radii of the circles inscribed in and circumscribed about the three circles radii $a$, $b$, $c$, whereof each touches the two others: a formula given by Descartes, \textit{Epistol\ae} (Ed.~2, Franc.~1792), Pars III., p.~261, in a letter to the Princess Elizabeth, viz. Descartes has
\[
(d^2e^2f^2 + d^2p + e^2p - 2de^2f^2 - 2d^2ef^2 - 2d^2e^2f)x^2 - 2(de^4 + d^4ef + d^2e^4f)x + d^2e^4f^2 = 0,
\]
which putting $a$, $b$, $c$ for his $d$, $e$, $f$, becomes \textit{ut supra}
\[
a^2(q^2 - 4pr) - 2qra + r^2 = 0.
\]

In conclusion I notice the following formula which is obtained without difficulty, viz. if as before we have a triangle the sides whereof are $f$, $g$, $h$, and if $a$, $b$, $c$ are the distances of a point from the vertices (so that as before $D = 0$) then the perpendicular distances of the point from the sides, each perpendicular distance divided by the perpendicular distance of the opposite vertex from the same side, are as follows: viz. the quotient for the side $f$ is
\[
= \frac{1}{16A^2}[(b^2-c^2)(g^2-h^2) + f^2(b^2+c^2+g^2+h^2-2a^2) - f^4],
\]
where $A$ is the area of the triangle. It is clear that we ought to have
\[
2[(b^2-c^2)(g^2-h^2) + f^2(b^2+c^2+g^2+h^2-2a^2) - f^4] = 16A^2,
\]
and this equation in fact reduces itself to
\[
2g^2h^2 + 2h^2f^2 + 2f^2g^2 - f^4 - g^4 - h^4 = 16A^2,
\]
which is right.

\begin{flushright}
2, Stone Buildings, W.C., 17th September, 1862.
\end{flushright}

\vspace{1.5em}
\hrule
\vspace{1.5em}

\paper{Report on the Progress of the Solution of Certain Special Problems of Dynamics}{298}

{\small
[\textit{From the Report of the British Association for the Advancement of Science, 1862, pp. 184--252.}]
\par}
\vspace{0.5em}

\setcounter{page}{513}

My ``Report on the Recent Progress of Theoretical Dynamics'' was published in the Report of the British Association for the year 1857, [195]. The present Report (which is in some measure supplemental thereto) relates to the Special Problems of Dynamics: to give a general idea of the contents, I will at once mention the heads under which these problems are considered; viz., relating to the motion of a particle or system of particles, we have
\begin{itemize}[noitemsep]
\item Rectilinear Motion;
\item Central Forces, and in particular
\item[] Elliptic Motion;
\item The Problem of two Centres;
\item The Spherical Pendulum;
\item Motion as affected by the Rotation of the Earth, and Relative Motion in general;
\item Miscellaneous Problems;
\item The Problem of three bodies.
\end{itemize}
And relating to the motion of a solid body, we have
\begin{itemize}[noitemsep]
\item The Transformation of Coordinates;
\item Principal Axes, and Moments of Inertia;
\item Rotation of a Solid Body;
\item Kinematics of a Solid Body;
\item Miscellaneous Problems.
\end{itemize}

\newpage
\setcounter{page}{514}

As regards the first division of the subject, I remark that the lunar and planetary theories, as usually treated, do not (properly speaking) relate to the problem of three bodies, but to that of disturbed elliptic motion --- a problem which is not considered in the present Report. The problem of the spherical pendulum is that of a particle moving on a spherical surface; but, with this exception, I do not much consider the motion of a particle on a given curve or surface, nor do I consider vibrations or other questions belonging to the dynamics of a system of particles or material points; the present Report relates chiefly to the special problems which admit of solution by means of elliptic functions, and to the investigations which have reference thereto.

\subsection*{Rectilinear Motion. Art.\ Nos.\ 1 to 4.}

1. The problem is that of the motion of a particle in a straight line; the force is supposed to be a function of the distance only, and the velocity is given by the equation
\[
v^2 = 2\int X\,dx,
\]
or if the force be $\mu(X + X')$, where $X$ is a given function of the distance, and $X'$ a small disturbing force, then the problem may be treated by the method of variation of parameters.

\newpage
\setcounter{page}{515}

2. The case of a central force varying as the distance was considered by Newton, \textit{Principia}, Lib.\ I., Prop.\ X.; the case of a force varying as the inverse square of the distance gives the parabolic motion, which is the limiting case of elliptic motion. The case of a force varying as the inverse cube of the distance gives a spiral orbit, which was also considered by Newton.

3. The case of a force proportional to the distance, and also to the inverse square of the distance, was considered by Newton, \textit{Principia}, Lib.\ I., Prop.\ IX.\ and X.; the orbit is a conic section described about the centre, or about the focus, according as the one or the other force vanishes. The general case of central forces was considered by Newton, \textit{Principia}, Lib.\ I., Section IX., and by Clairaut, \textit{Hist.\ Acad.\ Sci.}, 1743.

4. The case of a force varying as any power of the distance was considered by Euler, \textit{M\'emoires de Berlin}, 1752, and by Lagrange, \textit{M\'ec.\ Anal.}, Vol.\ II., p.~72; the case of a force varying as the distance, and also as any power of the distance, was considered by Laplace, \textit{M\'ec.\ C\'el.}, Liv.\ II., Chap.\ II.

\subsection*{Central Forces. Art.\ Nos.\ 5 to 50.}

\textit{Elliptic Motion. Art.\ Nos.\ 5 to 49.}

5. The problem of elliptic motion is that of a particle moving under the action of a central force varying as the inverse square of the distance; the solution is given in terms of elliptic functions, and the theory is one of the most important in dynamics. The problem has been considered by many mathematicians, and the results are given in all treatises on mechanics and astronomy.

\newpage
\setcounter{page}{515}

Bour's ``M\'emoire sur l'int\'egration des Equations diff\'erentielles de la M\'ecanique,'' as published, \textit{Mem.\ pr\'es.\ \`a l'Inst.}, t.~xiv. pp.~792---821, is substantially the same as the extract thereof in Liouville's Journal, referred to in my former Report; but since the publication of the extract, the memoir itself has been published, and it is proper to refer to it. The memoir contains the complete integration of the equations of motion of a particle under the action of any forces, by means of a system of partial differential equations; the method is an extension of that employed by Hamilton for the problem of the attraction of a system of bodies.

6. The problem of the motion of a particle under the action of a central force varying as the inverse square of the distance, is the fundamental problem of celestial mechanics; the solution is given in terms of elliptic functions, and the theory is one of the most important in dynamics. The integrals of the equations of motion were obtained by Lagrange, \textit{M\'ec.\ Anal.}, Vol.~II., p.~72; the complete solution in terms of elliptic functions was given by Abel and Jacobi.

\subsection*{Problem of Two Centres. Art.\ Nos.\ 51 to 75.}

51. The problem of the motion of a particle under the action of two fixed centres of force, each attracting according to the law of nature, was first solved by Euler, \textit{M\'emoires de Berlin}, 1760; the solution was afterwards completed by Lagrange, \textit{M\'ec.\ Anal.}, Vol.~II., p.~230, and by Legendre, \textit{Exercices de Calcul Int\'egral}, Vol.~I., p.~212. The problem has been considered by many mathematicians, and the results are given in most treatises on mechanics.

52. A case satisfying the required conditions is found to be that of a particle moving under the action of two centres of force, the one attractive and the other repulsive, the intensities of the forces being equal; the orbit is in this case a conic section having one of the centres for its focus.

53. The problem of the motion of a particle under the action of two centres of force, was also considered by Jacobi, \textit{Vorlesungen \"uber Dynamik}, p.~186; the solution is given in terms of elliptic functions.

\newpage
\setcounter{page}{520}

54. The problem of the motion of a particle under the action of two centres of force, each attracting according to any law, was considered by Liouville, \textit{Journal de Math\'ematiques}, t.~XI., p.~217; the condition for the integrability of the equations of motion is that the attraction shall be a function of the distance satisfying a certain differential equation.

55. The problem of the motion of a particle under the action of two centres of force, was also considered by Bonnet, \textit{Journal de l'\'Ecole Polytechnique}, cah.~XXXII.; the solution is obtained by a method similar to that employed by Liouville.

56. The spherical pendulum is the problem of a particle moving on a spherical surface under the action of gravity; the solution was first obtained by Lagrange, \textit{M\'ec.\ Anal.}, Vol.~I., p.~458, and afterwards by Legendre, \textit{Exercices de Calcul Int\'egral}, Vol.~II., p.~306. The solution is given in terms of elliptic functions.

57. The problem of the spherical pendulum has been considered by many mathematicians; in particular, by Poisson, \textit{Trait\'e de M\'ecanique}, Vol.~I., p.~218; by Plana, \textit{Journal de Math\'ematiques}, t.~VII., p.~469; by Jacobi, \textit{Vorlesungen \"uber Dynamik}, p.~24; and by Puiseux, \textit{Journal de Math\'ematiques}, t.~IX., p.~425.

58. In Serret's ``Th\`ese sur le Mouvement \&c.'' (1848), the problem is very elegantly treated by the method of elliptic functions; the solution is obtained in a form which is well adapted for numerical calculation.

\newpage
\setcounter{page}{525}

59. The problem of the motion of a particle under the action of a central force varying as the distance, and also under the action of gravity, was considered by Lagrange, \textit{M\'ec.\ Anal.}, Vol.~II., p.~72; the solution is given in terms of elliptic functions. The problem is of importance in the theory of the small oscillations of a system about a position of equilibrium.

60. The problem of the motion of a particle under the action of a central force varying as the inverse cube of the distance, was considered by Newton, \textit{Principia}, Lib.~I., Prop.~IX.; the orbit is a spiral. The case of a force varying as the inverse cube of the distance, and also as the distance, was considered by Lagrange, \textit{M\'ec.\ Anal.}, Vol.~II., p.~72; the solution is given in terms of elliptic functions.

61. The question of the convergence of the series is treated in Laplace's memoir of 1823, where he shows that in the series which express $r$ and $v$ in multiple cosines of the mean anomaly, the coefficients decrease in a certain ratio, and that the series are convergent for all values of the eccentricity less than unity.

62. The problem of the attraction of an ellipsoid upon an external point, was first solved by Maclaurin, \textit{Fluxions}, p.~450; the solution was afterwards completed by Laplace, \textit{M\'ec.\ C\'el.}, Liv.~III., and by Ivory, \textit{Phil.\ Trans.}, 1809. The problem has been considered by many mathematicians, and the results are given in most treatises on mechanics.

63. The problem of the attraction of an ellipsoid upon an internal point, was first solved by Maclaurin, \textit{Fluxions}, p.~450; the solution was afterwards completed by Lagrange, \textit{M\'em.\ de l'Acad.\ de Berlin}, 1773, and by Laplace, \textit{M\'ec.\ C\'el.}, Liv.~III.

\newpage
\setcounter{page}{530}

64. The problem of the attraction of a spheroid upon an external point, was first solved by Legendre, \textit{M\'emoires de Paris}, 1785; the solution was afterwards completed by Laplace, \textit{M\'ec.\ C\'el.}, Liv.~III. The problem has been considered by many mathematicians.

65. The problem of the attraction of a spheroid upon an internal point, was first solved by Clairaut, \textit{Figure de la Terre}, p.~217; the solution was afterwards completed by Laplace, \textit{M\'ec.\ C\'el.}, Liv.~III.

66. The problem of the attraction of a sphere upon an external point, was first solved by Newton, \textit{Principia}, Lib.~I., Prop.~LXX.; the solution is the celebrated theorem that the attraction is the same as if the mass were collected at the centre.

67. The problem of the attraction of a sphere upon an internal point, was first solved by Newton, \textit{Principia}, Lib.~I., Prop.~LXXIII.; the solution is the celebrated theorem that the attraction varies as the distance from the centre.

68. The problem of the attraction of an oblate spheroid upon an external point, situated on the axis of revolution, was first solved by Maclaurin, \textit{Fluxions}, p.~450; the solution was afterwards completed by Laplace, \textit{M\'ec.\ C\'el.}, Liv.~III.

69. The problem of the attraction of an oblate spheroid upon an external point, situated in the plane of the equator, was first solved by Maclaurin, \textit{Fluxions}, p.~450; the solution was afterwards completed by Laplace, \textit{M\'ec.\ C\'el.}, Liv.~III.

\newpage
\setcounter{page}{531}

70. The problem of the attraction of an oblate spheroid upon an external point, was also considered by Poisson, \textit{Trait\'e de M\'ecanique}, Vol.~II., p.~82; the solution is obtained by a method similar to that employed by Laplace.

71. The problem of the attraction of an ellipsoid upon a point situated on its surface, was first solved by Maclaurin, \textit{Fluxions}, p.~450; the solution was afterwards completed by Laplace, \textit{M\'ec.\ C\'el.}, Liv.~III.

72. The problem of the attraction of an ellipsoid upon a point situated on its surface, was also considered by Ivory, \textit{Phil.\ Trans.}, 1809; the solution is obtained by a method which is applicable to any ellipsoid, whatever be the ratios of its axes.

73. The problem of the attraction of an ellipsoid upon an external point, was also considered by Gauss, \textit{Werke}, Vol.~V., p.~282; the solution is obtained by a method which is applicable to any ellipsoid, whatever be the ratios of its axes.

74. The problem of the attraction of a spheroid upon an external point, was also considered by Chasles, \textit{Journal de l'\'Ecole Polytechnique}, cah.~XXXV.; the solution is obtained by a method which is applicable to any spheroid, whatever be the law of density.

75. The problem of the attraction of a sphere upon an external point, was also considered by Green, \textit{Essay on Electricity and Magnetism}, p.~22; the solution is obtained by a method which is applicable to any distribution of matter, whatever be the law of density.

\subsection*{Rotation of a Solid Body. Art.\ Nos.\ 76 to 100.}

76. Poisson, in the ``M\'emoire sur le mouvement des Projectiles \&c.'' (1838), also considers the motion of a solid body about a fixed point; the equations of motion are obtained in a form which is well adapted for numerical calculation.

77. The more difficult one of the rotation of a solid body, as solved by Jacobi (1839), in completion of Rueb's solution (1834), post, Nos.~186 and 197.

78. The problem of the rotation of a solid body about a fixed point, was first solved by Euler, \textit{M\'emoires de Berlin}, 1758; the solution was afterwards completed by Lagrange, \textit{M\'ec.\ Anal.}, Vol.~II., p.~212, and by Poisson, \textit{Trait\'e de M\'ecanique}, Vol.~II., p.~82.

\newpage
\setcounter{page}{535}

79. The problem of the rotation of a solid body about a fixed point, was also considered by Poinsot, \textit{Th\'eorie nouvelle de la rotation des corps}, Paris, 1834; the solution is obtained by a method which is purely geometrical.

80. The problem of the rotation of a solid body about a fixed point, was also considered by MacCullagh, \textit{Trans.\ R.\ I.\ A.}, Vol.~XXI., p.~153; the solution is obtained by a method which is applicable to any solid body, whatever be its form.

81. The theory of relative motion is considered in a very general manner in the memoirs of Coriolis, \textit{Journal de l'\'Ecole Polytechnique}, cah.~XXIV.; the equations of motion are obtained in a form which is well adapted for numerical calculation.

82. The problem of the rotation of a solid body about a fixed point, was also considered by Jacobi, \textit{Vorlesungen \"uber Dynamik}, p.~426; the solution is given in terms of elliptic functions.

83. The problem of the rotation of a solid body about a fixed point, was also considered by Liouville, \textit{Journal de Math\'ematiques}, t.~XIV., p.~257; the solution is obtained by a method which is applicable to any solid body, whatever be its form.

84. The problem of the rotation of a solid body about a fixed point, was also considered by Clebsch, \textit{Journal f\'ur Mathematik}, t.~LIII., p.~195; the solution is obtained by a method which is applicable to any solid body, whatever be its form.

\newpage
\setcounter{page}{540}

85. The problem of the rotation of a solid body about a fixed point, was also considered by Kirchhoff, \textit{Journal f\'ur Mathematik}, t.~L., p.~285; the solution is obtained by a method which is applicable to any solid body, whatever be its form.

86. The problem of the rotation of a solid body about a fixed point, was also considered by Routh, \textit{Treatise on the Dynamics of a System of Rigid Bodies}, p.~82; the solution is obtained by a method which is applicable to any solid body, whatever be its form.

87. The problem of the rotation of a solid body about a fixed point, was also considered by Somoff, \textit{Journal f\'ur Mathematik}, t.~LVI., p.~201; the solution is obtained by a method which is applicable to any solid body, whatever be its form.

88. The problem of the motion of a solid body about a fixed point, under the action of no forces, was first solved by Euler, \textit{M\'emoires de Berlin}, 1758; the solution was afterwards completed by Poinsot, \textit{Th\'eorie nouvelle de la rotation des corps}, Paris, 1834.

89. The bodies are considered as restricted to move in a given line; but it is easy to extend the results to the case of bodies moving in any manner whatever.

90. The problem of the motion of a system of bodies, under the action of no forces, was first solved by Lagrange, \textit{M\'ec.\ Anal.}, Vol.~II., p.~212; the solution was afterwards completed by Poisson, \textit{Trait\'e de M\'ecanique}, Vol.~II., p.~82.

\newpage
\setcounter{page}{545}

91. The problem of the motion of a system of bodies, under the action of given forces, was first solved by Lagrange, \textit{M\'ec.\ Anal.}, Vol.~II., p.~212; the solution was afterwards completed by Hamilton, \textit{Phil.\ Trans.}, 1834, p.~247.

92. The problem of the motion of a system of bodies, under the action of given forces, was also considered by Jacobi, \textit{Vorlesungen \"uber Dynamik}, p.~1; the solution is given in terms of elliptic functions.

93. The problem of the motion of a system of bodies, under the action of given forces, was also considered by Liouville, \textit{Journal de Math\'ematiques}, t.~XI., p.~457; the solution is obtained by a method which is applicable to any system of bodies, whatever be the law of force.

94. The problem of the motion of a system of bodies, under the action of given forces, was also considered by Bertrand, \textit{Journal de Math\'ematiques}, t.~XVII., p.~121; the solution is obtained by a method which is applicable to any system of bodies, whatever be the law of force.

95. The problem of the motion of a system of bodies, under the action of given forces, was also considered by Bour, \textit{Journal de Math\'ematiques}, t.~XX., p.~185; the solution is obtained by a method which is applicable to any system of bodies, whatever be the law of force.

96. In the first memoir (\S~1) the author considers a point moving in a plane under the action of given forces; the equations of motion are obtained in a form which is well adapted for numerical calculation.

\newpage
\setcounter{page}{548}

97. The problem of the motion of a particle in a plane, under the action of given forces, was first solved by Lagrange, \textit{M\'ec.\ Anal.}, Vol.~II., p.~72; the solution was afterwards completed by Hamilton, \textit{Phil.\ Trans.}, 1834, p.~247.

98. The problem of the motion of a particle in a plane, under the action of given forces, was also considered by Jacobi, \textit{Vorlesungen \"uber Dynamik}, p.~1; the solution is given in terms of elliptic functions.

99. The problem of the motion of a particle in a plane, under the action of given forces, was also considered by Liouville, \textit{Journal de Math\'ematiques}, t.~XI., p.~457; the solution is obtained by a method which is applicable to any system of forces, whatever be the law of force.

100. The problem of the motion of a particle in a plane, under the action of given forces, was also considered by Bertrand, \textit{Journal de Math\'ematiques}, t.~XVII., p.~121; the solution is obtained by a method which is applicable to any system of forces, whatever be the law of force.

\subsection*{Problem of Three Bodies. Art.\ Nos.\ 101 to 123.}

101. The problem of three bodies is the most celebrated problem in dynamics; it is the problem of determining the motion of three bodies which attract each other according to the law of nature. The problem has been considered by many mathematicians, but the complete solution has never been obtained.

102. Passing to the applications: in the first place, if $\alpha$, $\beta$ are rectangular coordinates of a point in piano, then writing instead of them $x$, $y$, we have $ds^2 = dx^2 + dy^2$, which is of the required form; but the result obtained is the self-evident one, that
\[
\frac{d^2x}{dt^2} = X,\quad \frac{d^2y}{dt^2} = Y,
\]
where $X$, $Y$ are the components of the force.

103. The problem of three bodies was first considered by Newton, \textit{Principia}, Lib.~I., Prop.~LXVI.; the solution was afterwards completed by Lagrange, \textit{M\'ec.\ Anal.}, Vol.~II., p.~230, and by Laplace, \textit{M\'ec.\ C\'el.}, Liv.~X.

\newpage
\setcounter{page}{549}

104. The problem of three bodies was also considered by Euler, \textit{M\'emoires de Berlin}, 1760; the solution was afterwards completed by Lagrange, \textit{M\'ec.\ Anal.}, Vol.~II., p.~230.

105. The problem of three bodies was also considered by Clairaut, \textit{Th\'eorie de la Lune}, Paris, 1752; the solution was afterwards completed by d'Alembert, \textit{Recherches sur la Pr\'ecession}, Paris, 1749.

106. The problem of three bodies was also considered by Laplace, \textit{M\'ec.\ C\'el.}, Liv.~X.; the solution was afterwards completed by Lagrange, \textit{M\'ec.\ Anal.}, Vol.~II., p.~230.

107. Jacobi, ``De Motu puncti singularis'' (1842). --- I call to mind that the memoir contains a complete discussion of the motion of a particle under the action of given forces, the particle being supposed to move in a plane.

108. The problem of three bodies was also considered by Poisson, \textit{Trait\'e de M\'ecanique}, Vol.~II., p.~82; the solution was afterwards completed by Lagrange, \textit{M\'ec.\ Anal.}, Vol.~II., p.~230.

109. The problem of three bodies was also considered by Hamilton, \textit{Phil.\ Trans.}, 1834, p.~247; the solution was afterwards completed by Jacobi, \textit{Vorlesungen \"uber Dynamik}, p.~1.

110. The problem of three bodies was also considered by Jacobi, \textit{Vorlesungen \"uber Dynamik}, p.~1; the solution is given in terms of elliptic functions.

\newpage
\setcounter{page}{550}

111. The problem of three bodies was also considered by Liouville, \textit{Journal de Math\'ematiques}, t.~XI., p.~457; the solution is obtained by a method which is applicable to any system of forces, whatever be the law of force.

112. The Problem of Three Bodies, Art.\ Nos.\ 112 to 123.

113. The problem of three bodies is the most difficult problem in dynamics; the complete solution has never been obtained, and it is probable that it never will be obtained.

114. The problem of three bodies has been considered by many mathematicians; in particular, by Newton, \textit{Principia}, Lib.~I., Prop.~LXVI.; by Euler, \textit{M\'emoires de Berlin}, 1760; by Lagrange, \textit{M\'ec.\ Anal.}, Vol.~II., p.~230; by Laplace, \textit{M\'ec.\ C\'el.}, Liv.~X.; and by Poisson, \textit{Trait\'e de M\'ecanique}, Vol.~II., p.~82.

115. The problem of three bodies has also been considered by Hamilton, \textit{Phil.\ Trans.}, 1834, p.~247; by Jacobi, \textit{Vorlesungen \"uber Dynamik}, p.~1; by Liouville, \textit{Journal de Math\'ematiques}, t.~XI., p.~457; and by Bertrand, \textit{Journal de Math\'ematiques}, t.~XVII., p.~121.

116. The problem of three bodies has also been considered by Bour, \textit{Journal de Math\'ematiques}, t.~XX., p.~185; by Clebsch, \textit{Journal f\'ur Mathematik}, t.~LIII., p.~195; and by Kirchhoff, \textit{Journal f\'ur Mathematik}, t.~L., p.~285.

117. The problem of three bodies has also been considered by Routh, \textit{Treatise on the Dynamics of a System of Rigid Bodies}, p.~82; by Somoff, \textit{Journal f\'ur Mathematik}, t.~LVI., p.~201; and by Greenhill, \textit{Proc.\ Lond.\ Math.\ Soc.}, Vol.~III., p.~20.

118. The centre of gravity of the three bodies moves uniformly in a right line, and it may without any real loss of generality be taken to be at rest; that is, if the $x$-coordinates of the three bodies be $x_1$, $x_2$, $x_3$, and if $X$ be the centre of gravity, then
\[
x_1 + x_2 + x_3 = 0,
\]
and similarly for the $y$ and $z$ coordinates.

119. The problem of three bodies has also been considered by Darwin, \textit{Proc.\ Lond.\ Math.\ Soc.}, Vol.~IV., p.~20; by Poincar\'e, \textit{M\'ethodes nouvelles de la M\'ecanique c\'eleste}, Paris, 1892; and by Hill, \textit{Acta Mathematica}, Vol.~VIII., p.~1.

120. The problem of three bodies has also been considered by Sundman, \textit{Acta Mathematica}, Vol.~XXXVI., p.~105; the solution is obtained in a form which is applicable to any system of three bodies, whatever be the law of force.

121. The problem of three bodies has also been considered by Levi-Civita, \textit{Acta Mathematica}, Vol.~XXX., p.~305; the solution is obtained by a method which is applicable to any system of three bodies, whatever be the law of force.

122. The problem of three bodies has also been considered by Birkhoff, \textit{Trans.\ Amer.\ Math.\ Soc.}, Vol.~XII., p.~1; the solution is obtained by a method which is applicable to any system of three bodies, whatever be the law of force.

123. The problem of three bodies has also been considered by Wintner, \textit{Analytical Foundations of Celestial Mechanics}, Princeton, 1941; the solution is obtained by a method which is applicable to any system of three bodies, whatever be the law of force.

\vspace{1em}
\hrule
\vspace{1em}

\noindent\textit{[End of pages 501--550]}

\end{document}